\documentclass[12pt]{article}

\usepackage[margin=1in]{geometry}
\usepackage[T1]{fontenc}
\usepackage[utf8]{inputenc}
\usepackage{newtxtext}
\usepackage{microtype}
\usepackage{setspace}
\usepackage{xurl}
\usepackage{hyperref}
\usepackage{fancyhdr}

% ---- Document metadata (edit here) ----
\newcommand{\PaperTitle}{From ``Problem Student'' to Contextual Learner: Deconstructing Bias Toward Neurodivergent and CLD Students Through Practicum Observation}
\newcommand{\AuthorName}{Piter Garc\'ia}
\newcommand{\Affiliation}{University of Rochester, Warner School of Education}
\newcommand{\CourseName}{EDU 442}
\newcommand{\InstructorName}{Dr. O'Neil-White}
\newcommand{\DueDate}{April 13, 2026}
\newcommand{\AssignmentName}{Deconstructing Bias \textemdash{} Reflective Paper}

\hypersetup{
  hidelinks,
  pdftitle={\PaperTitle},
  pdfauthor={\AuthorName},
  pdfsubject={\CourseName\space\textemdash{} \AssignmentName},
}

\doublespacing
\setlength{\parindent}{0.5in}
\setlength{\parskip}{0pt}
\setlength{\emergencystretch}{3em}

\pagestyle{fancy}
\fancyhf{}
\fancyhead[R]{\thepage}
\renewcommand{\headrulewidth}{0pt}
\setlength{\headheight}{15pt}

\newcommand{\apaheading}[1]{\begin{center}\textbf{#1}\end{center}}
\newcommand{\reference}[1]{\noindent\hangindent=0.5in\hangafter=1 #1\par}
\newcommand{\EvidenceSet}[1]{\textsc{E#1}}
\newcommand{\JournalArtifact}[1]{\textsc{J#1}}

\begin{document}

\begin{titlepage}
\thispagestyle{fancy}
\begin{center}
\vspace*{0.14\textheight}
{\Large\bfseries \PaperTitle\par}

\vspace{0.06\textheight}
{\AuthorName\par}
{\Affiliation\par}

\vspace{1.25\baselineskip}
{\AssignmentName\par}

\vspace{1.25\baselineskip}
{\CourseName\par}
{\InstructorName\par}
{\DueDate\par}
\vfill
\end{center}
\end{titlepage}

% The titlepage environment resets the page counter; force APA-style numbering.
\setcounter{page}{2}

\apaheading{\PaperTitle}

This paper examines a bias I have carried into classroom practice: the tendency to interpret neurodivergent and culturally/linguistically diverse (CLD) students through a compliance-and-deficit lens. In my practicum, that bias can show up as an automatic story: if a student is loud, impulsive, emotionally escalated, or ``not listening,'' they must be choosing to be disrespectful or disruptive. For this assignment, I analyze in-person classroom interactions from my elementary practicum placement during spring 2026 using time-stamped, verbatim observation records from lessons and mentor feedback conversations, alongside one course seminar discussion on linguistic justice. I label these materials as Evidence Sets \EvidenceSet{1}--\EvidenceSet{4} and catalog them in \path{artifacts/evidence_index.md}. I also include Journal Artifacts \JournalArtifact{1}--\JournalArtifact{3} (identity/journal writing and a course roundtable excerpt) to make my positionality and intended throughline auditable.

The specific bias under investigation is the ``problem student'' stereotype as it attaches to difference: difference in regulation (often read through disability/neurodivergence) and difference in language/discourse (often read through culture and linguistic background). My guiding question is: \textbf{How do my practicum observations of neurodivergent and CLD students---especially students positioned as ``problematic''---challenge the stereotype that disruption is a fixed trait, and what changes when adults respond relationally rather than punitively?} That question keeps me anchored in both the community (students in my placement classroom) and my own role as an emerging teacher.

In a course roundtable discussion, I named the interpretive shift I am trying to practice: in school settings we often interpret ``visible frustration, noise avoidance or emotional escalation'' as ``defiance rather than as evidence'' of confusion, stress, and unmet instructional needs (\JournalArtifact{3}). That statement is the paper's throughline.

My central claim is that what gets labeled ``behavior'' is often communication: a contextual response to unclear expectations, cognitive overload, linguistic mismatch, or unmet regulation needs. When the conditions shift---through relationship, co-regulation, and more linguistically just participation norms---students' engagement shifts, revealing competence and brilliance that deficit stories erase.

\apaheading{Autobiographical Grounding: Naming the Bias}

This paper is not written from a neutral position. I am writing as a neurodivergent Dominican-American student teacher who is simultaneously learning classroom management and learning to question the interpretive habits I bring into that work. The bias I am examining is not only something I have observed in others; it is a habit of interpretation that I recognize in my own instincts and that I need to examine honestly. I carry stories from being positioned as ``too much'' or ``too different'' in school spaces, and those stories sit beside my current authority as the adult in the room.

Talusan (2022) describes identity-conscious practice as the ongoing work of noticing which assumptions we carry and how they shape what we see in students---not a one-time confrontation but a repeatable discipline. That framing captures my starting point. The stereotype I am interrogating---that certain students are naturally disruptive, lazy, or unmotivated---is not a belief I endorse on purpose; it is a reflex shaped by schooling environments in which compliance read as competence and visible struggle read as character deficit.

In an EDU 442 self-portrait earlier this term (\JournalArtifact{1}), I described moving through institutions that often treat culture, language, disability, and gender variance as deficits, even when they are sources of insight and resilience. I used a puzzle metaphor to reject a fragmented view of self: Dominican heritage, first-generation class position, ADHD, chronic illness, and culturally/linguistically diverse life interlock rather than compete. That metaphor matters here because the ``problem student'' stereotype is a story about which ways of speaking, moving, and processing are treated as legitimate.

This history makes me protective of neurodivergent and CLD students, but it does not automatically inoculate me against bias. Under time pressure or public scrutiny, it is easy to reach for control. In a journal letter to my future teacher-self (\JournalArtifact{2}), I worried about systems that ``say equity'' but still reward compliance. I also wrote that I wanted to avoid ``playing school'' by pretending that my own safety and regulation do not matter; instead, I need to make plans for when I feel triggered so that I do not pass that stress back to students. This paper is an attempt to practice a different reflex: to pause, ask what a behavior is communicating, and design conditions where more students can be seen as fully human learners.

\apaheading{Observation Context}

The evidence base for this paper comes from my elementary school practicum placement in spring 2026 and one EDU 442 seminar discussion. I cite these materials as Evidence Sets \EvidenceSet{1}--\EvidenceSet{4} (see \path{artifacts/evidence_index.md} for source files and excerpt line references). Altogether, the data set covers four distinct observations (approximately two hours of recorded classroom interaction and a ninety-minute seminar), which allows me to see patterns and not just one dramatic moment.

The primary focal setting is Evidence Set \EvidenceSet{1}: a grade 5 media/technology lesson (March 3, 2026) in which students learned to create short GIF videos on their Chromebooks. During this lesson, a focal student (pseudonym \textit{Student M}) was repeatedly framed through a behavioral deficit narrative (``screaming,'' ``yelling''), while adults alternated between public compliance demands and more relational de-escalation moves (for example, offering responsibility and negotiated incentives).

I also draw on Evidence Set \EvidenceSet{2}: a grade 2 robotics/routines block (February 26, 2026) in which an adult explicitly named one student's behavior as ``ruining this for your whole class'' and ``doing weird stuff,'' and then later repositioned that student as competent and helpful (``You're smart, you can help them'').

Evidence Set \EvidenceSet{3} is a mentor feedback conversation (April 8, 2026) in which my mentor discussed how relational connection and patience changed what was possible for \textit{Student M} across the school day. In this conversation we explicitly named my own triggers and the pressures I feel to ``have control'' in front of other adults.

Finally, I use Evidence Set \EvidenceSet{4}: an EDU 442 seminar discussion on linguistic justice (March 9, 2026) in which we named how dominant language and discourse norms shape what counts as ``respect'' and ``participation'' in schools. I use that discussion to think about how CLD students can be misread when schools treat silence, speed, and standardized speech as the only legitimate way to learn.

Across these observations I participated as a student teacher, not as a neutral outside researcher: I was co-teaching, circulating, and sometimes debriefing my own moves with a mentor. Before writing about these moments, I obtained my mentor's consent, removed all identifying information, and created pseudonyms for students. I also attend to my own safety and well-being by noticing when particular behaviors or adult comments spike my stress and debriefing those moments with peers and faculty, instead of letting them silently harden into stereotypes.

Ethically, all student identifiers are removed. Student names in quoted excerpts are replaced with pseudonyms in square brackets. Ellipses (\ldots{}) indicate omitted words for brevity; quoted phrasing otherwise preserves the original wording and disfluencies of spoken talk. Because Evidence Sets \EvidenceSet{1}--\EvidenceSet{4} are unpublished practicum/course observation materials, they are cited in-text but not included in the References list.

\apaheading{Myth vs. Reality: Core Analysis}

\textbf{The myth.} Neurodivergent and culturally/linguistically diverse (CLD) students become ``problem students'' because they are naturally disruptive, disrespectful, or unwilling to learn. In practice, this myth is sustained by two linked assumptions: (1) that silence and compliance equal learning, and (2) that deviation from dominant behavioral and linguistic norms signals character deficit.

DisCrit (Annamma et al., 2013) and Flynn's (2024) affirmative non-tragedy model help me name why this is not a neutral mistake: deficit framings of difference become reasons to control rather than to support, shaping students' trajectories. Linguistic bias is part of the same apparatus. In our linguistic justice seminar (\EvidenceSet{4}), we named how dominant language and participation norms can harm CLD students when schools treat silence, speed, and standardized speech as the only legitimate ways to learn.

\textbf{The reality.} Across my practicum observation records (\EvidenceSet{1}--\EvidenceSet{3}), the ``problem student'' category is produced moment-by-moment through adult talk, task design, and relationship. In the grade 5 media lesson (\EvidenceSet{1}), adults sometimes framed the student's body and voice as the problem to be suppressed through public compliance demands.

Yet the same record captures a different set of moves that changed what was possible. Another adult invited \textit{Student M} into a purposeful helper role and reinforced specific success instead of shame; later, that adult narrated a shift in perception: ``I'm actually so surprised how how smart he is'' (\EvidenceSet{1}). The student's competence was not newly created; it was newly \textit{recognized} once the conditions supported regulation and re-entry.

A mentor feedback conversation makes this relational frame explicit. My mentor described how connection and patience changed the student's school day (``He's used to people just yelling at him all day'') and I named triggers and control rather than willfulness as a lens for understanding escalation (\EvidenceSet{3}). The linguistic justice discussion (\EvidenceSet{4}) extends the same insight to CLD students: when schools equate learning with silence and standardized speech, difference is easily misread as disrespect.

\apaheading{Assets and Funds of Knowledge}

An asset-based reading of the same observations reveals what the ``problem student'' story erases. Gonzalez, Moll, and Amanti (2005) argue that learners bring funds of knowledge from home and community that schools often fail to recognize. In my practicum records, those assets surfaced when adults shifted from labeling to invitation: in \EvidenceSet{1}, \textit{Student M} could contribute as a helper; in \EvidenceSet{2}, an adult explicitly repositioned a student as capable and needed: ``You're smart, you can help them'' (\EvidenceSet{2}).

Evidence Set \EvidenceSet{2} also contains language that centers cognitive diversity rather than pathologizing it: ``There's never just one way to solve a problem when you're coding. There's always multiple ways to figure it out, and the way that you guys think differently is a strength'' (\EvidenceSet{2}). This matters because it names neurodivergent brilliance as real and worth organizing learning around.

For CLD students, assets are also linguistic. In our linguistic justice seminar (\EvidenceSet{4}), we discussed inviting students to explain ideas ``in their own language'' as an equity move in computer science. Paris and Alim's (2014) culturally sustaining pedagogy helps me see why this matters: when schooling treats only one language variety and one participation style as legitimate, it erases communicative strength. hooks (1994) describes this as engaged pedagogy---staying curious about the whole person underneath the behavior and building conditions where students can participate without being erased.

These moments also changed what I learned from the group itself. Watching students lean over each other's screens to debug code, negotiate rules for turn-taking, or translate instructions for peers reminded me that community is built laterally, not just from the front of the room. The group taught me that there are already practices of mutual aid and sense-making in the classroom that I can either amplify or shut down, depending on how I respond.

\apaheading{Implications for Future Practice}

These observations change what I want to do before I assign a label to a student. Talusan (2022) argues that identity-conscious practice is built from repeatable habits---not grand gestures, but daily routines of noticing and questioning. My commitments going forward are habits I can rehearse in real time.

First, I want to treat escalation as information, not defiance. When I notice a student yelling, shut down, or talking out of turn, my first questions should be contextual: what is the student trying to communicate, what is the trigger, and what support would make participation possible in the next two minutes? This is aligned with Flynn's (2024) affirmative orientation: difficulty is real, but the response should be support and dignity rather than tragedy and control.

Second, I want to lead with relationship and co-regulation. In \EvidenceSet{1} and \EvidenceSet{3}, purposeful roles, patience, and competence-talk supported re-entry into learning; those moves are part of instruction, not optional extras. Thoughtfully connecting with this group means building predictable routines, checking in on how noise and pace land in their bodies, and noticing who is consistently positioned on the margins of talk so that I can redesign the task.

Third, I want to practice linguistic justice as classroom management. If schools reward ``respect'' primarily as silence, speed, and standardized speech, then CLD and neurodivergent students are set up to be misread (\EvidenceSet{4}; Paris \& Alim, 2014). As I wrote to my future teacher-self (\JournalArtifact{2}), the fear is not structure itself---it is structure used to purchase adult comfort through compliance. The goal is participation norms that are taught explicitly and designed for learning, while still being firm and humane. That includes assessing when a requested behavior genuinely protects safety and learning and when it mainly protects adult comfort.

Across all of this, I want to audit my language and my interpretations (Annamma et al., 2013; Nieto, 2004): am I labeling a student, or describing conditions and supports? hooks (1994) keeps the ethical center in view: scaffold before I shame, narrate competence before deficit, and design participation structures that make room for more than one way of being a learner.

\apaheading{Conclusion}

I began this paper with a bias that is easy to carry invisibly into classroom practice: the reflex to treat behavioral and communicative difference as evidence of fixed student character. Practicum observation gave me evidence to interrupt that reflex.

Across the evidence sets, I saw how quickly a student can be produced as a ``problem'' through adult interpretation---and how quickly that category can loosen when adults respond relationally. When adults shifted from public control to co-regulation and competence-talk, participation became possible again (\EvidenceSet{1}--\EvidenceSet{3}). The linguistic justice lens sharpened the stakes for CLD students (\EvidenceSet{4}): if schools treat silence, speed, and standardized speech as the only legitimate ways to learn, difference becomes punishable. Going forward, my commitment is to make the interpretive pause into pedagogy: to read behavior as communication, design conditions that widen participation, and use structure to support learning rather than to purchase compliance. In my own words, I want it to be less likely that another CLD, neurodivergent kid slips through cracks that were never accidents (\JournalArtifact{2}).

\newpage
\apaheading{References}

\reference{Annamma, S. A., Connor, D., \& Ferri, B. (2013). Dis/ability critical race studies (DisCrit): Theorizing at the intersections of race and dis/ability. \textit{Race Ethnicity and Education, 16}(1), 1--31.}

\reference{Flynn, S. (2024). Critical disability studies and the affirmative non-tragedy model: Presenting a theoretical frame for disability and child protection. \textit{Disability \& Society, 39}(2), 269--290. \url{https://doi.org/10.1080/09687599.2022.2070061}}

\reference{Gonzalez, N., Moll, L. C., \& Amanti, C. (Eds.). (2005). \textit{Funds of knowledge: Theorizing practices in households, communities, and classrooms.} Lawrence Erlbaum.}

\reference{hooks, b. (1994). \textit{Teaching to transgress: Education as the practice of freedom.} Routledge.}

\reference{Nieto, S. (2004). \textit{Affirming diversity: The sociopolitical context of multicultural education} (4th ed.). Allyn \& Bacon.}

\reference{Paris, D., \& Alim, H. S. (2014). What are we seeking to sustain through culturally sustaining pedagogy? \textit{Harvard Educational Review, 84}(1), 85--100.}

\reference{Talusan, L. (2022). \textit{The identity-conscious educator: Building habits and skills for a more inclusive school.} Solution Tree Press.}

\end{document}
